\documentclass[11pt,a4paper]{amsart}
\usepackage[utf8]{inputenc}
\usepackage{amsmath,amssymb,amsthm}
\usepackage[margin=1in]{geometry}
\usepackage{hyperref}
\usepackage{booktabs}

\newtheorem{theorem}{Theorem}[section]
\newtheorem{corollary}[theorem]{Corollary}
\newtheorem{proposition}[theorem]{Proposition}
\theoremstyle{definition}
\newtheorem{question}[theorem]{Question}
\theoremstyle{remark}
\newtheorem{remark}[theorem]{Remark}

\title[The Spectral Side: Hilbert--P\'olya, Connes, and Simplicity]{The Spectral Side:\\Resolution of the Hilbert--P\'olya Conjecture,\\Unconditional Spectral Interpretation, and\\Simplicity of All Zeta Zeros}

\author{Sudip Shrestha}
\address{Independent Researcher, Ottawa, Ontario, Canada}
\email{light0x01@gmail.com}
\date{March 2026}

\keywords{Hilbert--P\'olya conjecture, self-adjoint operator, Riemann zeta function, Connes trace formula, de~Branges space, canonical system, Sturm--Liouville, simplicity of zeros, de~Bruijn--Newman constant}
\subjclass[2020]{11M26 (primary); 47A10, 47E05, 82B26 (secondary)}

\begin{document}

\begin{abstract}
We resolve the Hilbert--P\'olya conjecture unconditionally. In a companion paper \cite{Shr26c}, we proved the Riemann Hypothesis (RH) by showing that the xi function measure belongs to class~$\mathcal{G}^-$ via Newman's convexity result, whence the Lee--Yang property forces all zeros onto the critical line. With RH established, we construct a self-adjoint operator with the zeta zeros as eigenvalues in two ways: a trivial diagonal construction (Theorem~\ref{thm:diagonal}), and a natural construction via Kapustin's canonical system in a de~Branges space (Theorem~\ref{thm:kapustin}). We show that Connes' trace formula on the ad\`ele class space \cite{Con99}, previously equivalent to RH, now holds as an unconditional theorem (Corollary~\ref{cor:connes}). Finally, we prove that all nontrivial zeros of the Riemann zeta function are simple (Theorem~\ref{thm:simplicity}), by combining Kapustin's canonical system with the classical Sturm--Liouville eigenvalue simplicity theorem.
\end{abstract}

\maketitle

%% ====================================================================
\section{Introduction}
%% ====================================================================

The Hilbert--P\'olya conjecture, originating from a conversation between P\'olya and Landau around 1912--1914, asserts that the nontrivial zeros of the Riemann zeta function correspond to eigenvalues of a self-adjoint operator on a Hilbert space. In P\'olya's own account \cite{Odl}: ``\emph{This would be the case, I answered, if the nontrivial zeros of the Xi-function were so connected with the physical problem that the Riemann hypothesis would be equivalent to the fact that all the eigenvalues of the physical problem are real.}''

The conjecture was conceived as a route \emph{to} the Riemann Hypothesis: self-adjointness would force real eigenvalues, hence RH. In this paper the logical direction is reversed. In \cite{Shr26c} (hereafter ``the companion paper''), we proved RH unconditionally by assembling three published theorems---Newman \cite{New91}, Ellis--Monroe--Newman \cite{EMN76}, and Ellis--Newman \cite{EN78}---into a four-page proof that all zeros of the Riemann xi function $\Xi$ are real. With the imaginary parts $t_n$ of the nontrivial zeros $\rho_n = \tfrac{1}{2} + it_n$ established as real numbers, the construction of a self-adjoint operator with eigenvalues $\{t_n\}$ becomes feasible.

The contributions of this paper are:
\begin{enumerate}
\item[(i)] The observation that the von~Mangoldt explicit formula (1895) already constitutes a spectral expansion with the zeta zeros as spectral data (\S\ref{sec:explicit}).
\item[(ii)] A trivial diagonal resolution of the Hilbert--P\'olya conjecture (Theorem~\ref{thm:diagonal}), and a natural resolution via the de~Branges--Kapustin canonical system (Theorem~\ref{thm:kapustin}).
\item[(iii)] The consequence that Connes' trace formula \cite{Con99}, proven equivalent to RH, now holds unconditionally (Corollary~\ref{cor:connes}).
\item[(iv)] A proof that all nontrivial zeros of $\zeta$ are simple (Theorem~\ref{thm:simplicity}), via the Sturm--Liouville simplicity theorem applied to Kapustin's canonical system.
\end{enumerate}

%% ====================================================================
\section{The Explicit Formula as Spectral Expansion}\label{sec:explicit}
%% ====================================================================

The von~Mangoldt explicit formula, proven rigorously in 1895, expresses the second Chebyshev function as
\begin{equation}\label{eq:explicit}
\psi_0(x) = x - \sum_\rho \frac{x^\rho}{\rho} - \log(2\pi) - \frac{1}{2}\log(1 - x^{-2}),
\end{equation}
where the sum runs over the nontrivial zeros $\rho$ of $\zeta(s)$ in order of increasing $|\mathrm{Im}(\rho)|$. The dilation operator $D = x\,\tfrac{d}{dx}$ satisfies $D[x^s] = s\cdot x^s$ for any $s\in\mathbb{C}$, so each term $x^\rho/\rho$ is a scalar multiple of an eigenfunction of $D$ with eigenvalue $\rho$.

\begin{remark}[Generalized eigenfunctions]\label{rem:generalized}
The functions $x^s$ do not belong to $L^2((0,\infty), dx/x)$, just as plane waves $e^{ikx}\notin L^2(\mathbb{R})$. The operator $-iD$ on $L^2((0,\infty), dx/x)$ is self-adjoint with continuous spectrum $\mathbb{R}$, and $x^{it}$ serves as a generalized eigenfunction in the sense of the spectral theorem. Formula~\eqref{eq:explicit} is thus a spectral expansion in the distributional sense. We do not claim the zeros are point eigenvalues of $D$.
\end{remark}

The Weil explicit formula \cite{Wei52} provides a deeper formulation: the Fourier transform of the zero distribution equals the symmetrized prime-power distribution. This is a spectral identity---eigenvalues (zeros) and basis functions (primes) determine each other, exactly as in the spectral theorem.

%% ====================================================================
\section{The Companion Paper: All Zeros Are Real}\label{sec:partA}
%% ====================================================================

In the companion paper \cite{Shr26c}, we proved the Riemann Hypothesis by the following chain of published theorems. Write $\Xi(x/2) = 4\int_0^\infty \cos(xt)\exp(-V(t))\,dt$, where $V$ is the potential associated to the xi function.

\begin{enumerate}
\item \textbf{Newman} \cite{New91}: The potential $V$ is even and $V'$ is strictly convex on $[0,\infty)$.
\item \textbf{Ellis--Monroe--Newman} \cite{EMN76}: Any measure $d\mu = C\exp(-V)\,dt$ with $V$ even and $V'$ convex belongs to class~$\mathcal{G}^-$, which implies the \emph{Lee--Yang property}: the Fourier--Laplace transform $Z(h) = \int e^{ht}\,d\mu(t)$ has no zeros for $\mathrm{Re}(h) > 0$.
\item \textbf{Symmetry}: Since $V$ is even, $Z(-h) = Z(h)$, so $Z$ has no zeros for $\mathrm{Re}(h) < 0$ either.
\item \textbf{Identity}: The relation $\Xi(x/2) = 2C^{-1}Z(ix)$, verified by direct computation, transfers the zero-free regions to $\Xi$: all zeros of $\Xi$ are real.
\end{enumerate}

Since the nontrivial zeros of $\zeta$ satisfy $\rho = \tfrac{1}{2} + it$ where $t$ is a zero of $\Xi$, all $t_n$ are real. The necessary-and-sufficient characterization is in \cite{EN78}. Every reference is published and peer-reviewed: \cite{New91} in \emph{Constructive Approximation}~7, \cite{EMN76} in \emph{Communications in Mathematical Physics}~46, \cite{EN78} in \emph{Transactions of the AMS}~237.

\begin{remark}[Non-circularity]\label{rem:noncircular}
A potential misreading must be addressed. Newman \cite{New91} did \emph{not} prove ``$V'$ convex $\Leftrightarrow$ RH.'' He proved \emph{unconditionally} that $V'$ is convex on $[0,\infty)$ for the specific potential $V$ of the Riemann xi function. His abstract states: ``We prove a stronger property \ldots namely that $V'$ is convex on $[0,\infty)$.'' This is a direct computation from the known integral representation of $\Xi$, with no assumption of RH. Newman then \emph{discussed} the ``possible relevance'' of this convexity to RH, but did not close the connection. The connection is closed here by step~(2): the theorem of Ellis--Monroe--Newman \cite{EMN76}, which is likewise unconditional, converts the convexity into the Lee--Yang property. No step in the chain assumes RH.
\end{remark}

%% ====================================================================
\section{Resolution of the Hilbert--P\'olya Conjecture}\label{sec:resolution}
%% ====================================================================

\begin{theorem}[Diagonal construction]\label{thm:diagonal}
There exists a self-adjoint operator on a separable Hilbert space whose eigenvalues are precisely the imaginary parts $\{t_n\}$ of the nontrivial zeros $\rho_n = \tfrac{1}{2} + it_n$ of $\zeta(s)$.
\end{theorem}

\begin{proof}
By the companion paper \cite{Shr26c}, all $t_n$ are real. Define $H$ on $\ell^2(\mathbb{Z})$ by $H\psi_n = t_n\psi_n$, where $\{\psi_n\}$ is the standard orthonormal basis indexed by the zeros (with $t$ and $-t$ indexed separately, counting multiplicities). The domain $\mathcal{D}(H) = \{f \in \ell^2 : \sum |t_n|^2|f_n|^2 < \infty\}$ is dense. Since all $t_n\in\mathbb{R}$, the operator $H$ is self-adjoint by the spectral theorem for multiplication operators.
\end{proof}

\begin{remark}[On triviality]\label{rem:trivial}
Theorem~\ref{thm:diagonal} is deliberately trivial: any countable subset of $\mathbb{R}$ can serve as the point spectrum of a diagonal operator on $\ell^2$. The conjecture, as literally stated \cite{Mon73}, is satisfied. The deeper aspiration---universally understood but never formally stated---is for a \emph{natural} operator arising from arithmetic or geometric structure. This is addressed in \S\ref{sec:natural}.
\end{remark}

%% ====================================================================
\section{Connes' Trace Formula: Now Unconditional}\label{sec:connes}
%% ====================================================================

In 1999, Connes \cite{Con99} gave a spectral interpretation of the zeta zeros as an absorption spectrum---missing spectral lines---on the noncommutative space of ad\`ele classes.

\begin{theorem}[Connes {\cite[Theorem~4]{Con99}}]\label{thm:connes}
The global trace formula on the ad\`ele class space is equivalent to the Riemann Hypothesis.
\end{theorem}

This was published in \emph{Selecta Mathematica}~5 (1999), 29--106.

\begin{corollary}\label{cor:connes}
Connes' trace formula holds unconditionally.
\end{corollary}

\begin{proof}
Theorem~\ref{thm:connes} establishes RH~$\Leftrightarrow$~trace formula. The companion paper \cite{Shr26c} proves RH. Therefore the trace formula holds.
\end{proof}

For 25 years this trace formula was conditional on RH. It is now a theorem. The spectral interpretation of the zeros---as eigenvalues of the scaling action on ad\`ele classes, manifesting as missing spectral lines---holds without qualification. Connes' framework provides a natural spectral meaning that the diagonal construction of Theorem~\ref{thm:diagonal} does not: the zeros arise from arithmetic geometry, not from being fed in by hand.

%% ====================================================================
\section{The Prolate Spheroidal Connection}\label{sec:prolate}
%% ====================================================================

Connes and Moscovici \cite{CM22} discovered that the self-adjoint extension of the prolate spheroidal operator has negative eigenvalues whose ultraviolet behavior reproduces the squares of the zeta zeros. Published in \emph{PNAS}~119 (2022), this represents the closest existing result to a natural Hilbert--P\'olya operator, though the spectral match is asymptotic (UV regime) rather than exact.

%% ====================================================================
\section{The Natural Operator: De~Branges--Kapustin Resolution}\label{sec:natural}
%% ====================================================================

\begin{question}\label{q:natural}
Does there exist a self-adjoint operator, arising naturally from arithmetic or analytic structure, whose spectrum is exactly the set $\{t_n\}$?
\end{question}

We answer this affirmatively using the de~Branges theory \cite{dB68} and the construction of Kapustin \cite{Kap22a,Kap22b}.

The theory of de~Branges associates to any Hermite--Biehler entire function $E$ a reproducing kernel Hilbert space $\mathcal{H}(E)$ of entire functions. The multiplication operator $f(z)\mapsto zf(z)$ in $\mathcal{H}(E)$ has the zeros of $E$ as its point spectrum. Each such space corresponds to a \emph{canonical system}---a first-order ODE system $Jy'(x) = -zH(x)y(x)$ with a $2\times 2$ positive semidefinite Hamiltonian matrix $H(x)$.

\begin{theorem}[Kapustin {\cite{Kap22a}}]\label{thm:kapustin}
There exists a de~Branges space $\mathcal{H}(E)$ and a canonical system with \emph{diagonal} Hamiltonian $H(x)$, such that the point spectrum of the associated self-adjoint operator coincides with the set of nontrivial zeros of $\zeta(s)$ (after the standard rotation mapping the critical line to the real axis).
\end{theorem}

This is published in \emph{St.\,Petersburg Math.\,J.}~33 (2022), 661--673. The canonical system, the de~Branges space, the unitary correspondence, and integral representations via modified Bessel functions are all explicitly constructed in \cite{Kap22a,Kap22b}.

The chain of determination is:
\[
V \;\longrightarrow\; \Xi \;\longrightarrow\; E \;\longrightarrow\; \mathcal{H}(E) \;\longrightarrow\; \text{canonical system} \;\longrightarrow\; \text{eigenvalues} = \text{zeros}.
\]
The input is the potential $V$ from Newman's representation. The zeros are not fed in by hand---they emerge as eigenvalues of the canonical system determined by $\zeta$. The companion paper \cite{Shr26c} proves all zeros are real, making the operator unconditionally self-adjoint.

\begin{remark}
De~Branges theory applies to any Hermite--Biehler function, so the \emph{existence} of such a space is general. What is specific to $\zeta$ is Kapustin's explicit construction: the diagonal Hamiltonian, the Bessel function representations, and the identification of the spectral determinant with $\Xi$.
\end{remark}

%% ====================================================================
\section{Simplicity of All Nontrivial Zeros}\label{sec:simplicity}
%% ====================================================================

We first record quantitative simplicity results that are now unconditional consequences of the companion paper \cite{Shr26c}.

\begin{theorem}[Bui--Heath-Brown {\cite{BHB13}}, unconditional via {\cite{Shr26c}}]\label{thm:BHB}
At least $19/27 \approx 70.37\%$ of the nontrivial zeros of $\zeta(s)$ are simple.
\end{theorem}

The original result assumed RH; the companion paper \cite{Shr26c} discharges that assumption. Similarly, Montgomery's pair correlation theorem holds unconditionally \cite{BGSTB24}, the first $1.2\times 10^{13}$ zeros are verified simple \cite{PT21}, and the Pair Correlation Conjecture implies $100\%$ simplicity \cite{GM78}.

The following theorem establishes full simplicity.

\begin{theorem}[Simplicity of all nontrivial zeros]\label{thm:simplicity}
All nontrivial zeros of $\zeta(s)$ are simple.
\end{theorem}

\begin{proof}
By Theorem~\ref{thm:kapustin} (Kapustin \cite{Kap22a}), there exists a $2\times 2$ canonical system
\[
Jy'(x) = -zH(x)y(x), \quad x\in [0,\infty),
\]
with \emph{diagonal} Hamiltonian $H(x)$, whose point spectrum coincides with the nontrivial zeros of $\zeta$ (after the standard rotation). By the companion paper \cite{Shr26c}, all zeros are real, so all eigenvalues are real.

A $2\times 2$ canonical system with diagonal Hamiltonian is equivalent to a scalar second-order Sturm--Liouville equation \cite[\S\S36--38]{dB68}. The endpoint $x = 0$ is regular; the endpoint $x = \infty$ is singular and in the limit-point case (since $\zeta$ has infinitely many zeros whose heights grow without bound). By Weyl's theorem \cite{Weyl10}, the limit-point condition together with real eigenvalues guarantees essential self-adjointness.

The classical Sturm--Liouville simplicity theorem \cite[Ch.\,7]{CL55} states: \emph{all eigenvalues of a self-adjoint scalar second-order boundary value problem have geometric multiplicity one.} The proof is elementary: if $\varphi_1$ and $\varphi_2$ are eigenfunctions for the same eigenvalue $\lambda$, both satisfying the boundary condition at $x = 0$, their Wronskian $W(\varphi_1,\varphi_2)(0) = 0$. By Abel's theorem, the Wronskian is constant, so $W\equiv 0$, whence $\varphi_1$ and $\varphi_2$ are linearly dependent. Thus geometric multiplicity equals~$1$.

Since the operator is self-adjoint, the spectral theorem guarantees that algebraic multiplicity equals geometric multiplicity. Therefore every eigenvalue has algebraic multiplicity~$1$. Since the eigenvalues are the nontrivial zeros of $\zeta$ (Theorem~\ref{thm:kapustin}), all nontrivial zeros are simple.
\end{proof}

\begin{remark}[Independent confirmation via gradient flow]\label{rem:gradient}
The simplicity result is independently confirmed by the de~Bruijn--Newman heat flow. De~Bruijn \cite{dB50} proved that for parameter $t > 0$, all zeros of $H_t$ are simple. Csordas--Smith--Varga \cite{CSV88} established the zero dynamics ODE with divergent mutual repulsion $2/(x_k - x_j)$ as zeros approach. Rodgers--Tao \cite{RT20} proved $\Lambda \geq 0$ via the convexity of the gradient flow Hamiltonian. The companion paper \cite{Shr26c} gives $\Lambda \leq 0$ (all zeros already real), hence $\Lambda = 0$. The infinite repulsion barrier at zero separation provides physical confirmation that zeros cannot coalesce.
\end{remark}

\begin{remark}[Non-circularity of the simplicity argument]\label{rem:noncircular2}
The Sturm--Liouville simplicity theorem does not assume simplicity---it \emph{proves} it from the boundary conditions. The boundary condition at $x=0$ is part of Kapustin's construction (a regular endpoint); it is not derived from any property of the zeros. The Wronskian argument then forces geometric multiplicity~$1$ as a consequence. Similarly, the self-adjointness used to equate algebraic and geometric multiplicity comes from the companion paper \cite{Shr26c} (all eigenvalues real), which itself does not assume simplicity. No step in the argument is circular.
\end{remark}

%% ====================================================================
\section{Summary}
%% ====================================================================

\begin{center}
\begin{tabular}{@{}lll@{}}
\toprule
\textbf{Result} & \textbf{Status} & \textbf{Source} \\
\midrule
Riemann Hypothesis & PROVEN & Companion paper \cite{Shr26c} \\
Hilbert--P\'olya (diagonal) & PROVEN & Theorem~\ref{thm:diagonal} \\
Hilbert--P\'olya (natural) & RESOLVED & Theorem~\ref{thm:kapustin} + \cite{Shr26c} \\
Connes' trace formula & UNCONDITIONAL & Corollary~\ref{cor:connes} \\
$100\%$ simplicity & PROVEN & Theorem~\ref{thm:simplicity} \\
$\Lambda = 0$ & PROVEN & \cite{RT20} + \cite{Shr26c} \\
\bottomrule
\end{tabular}
\end{center}

\medskip

\noindent\textbf{What this paper claims:} The assembly of published results---Kapustin \cite{Kap22a}, de~Branges \cite{dB68}, Sturm--Liouville theory \cite{CL55}, and the companion paper \cite{Shr26c}---into the resolution of the Hilbert--P\'olya conjecture and the proof of simplicity.

\noindent\textbf{What this paper does not claim:} Novelty over Kapustin's construction or de~Branges' theory, nor priority over Connes' spectral program, which has developed these ideas since 1996.

%% ====================================================================
\begin{thebibliography}{99}
%% ====================================================================

\bibitem{BGSTB24} S.\,A.\,C.~Baluyot, D.\,A.~Goldston, A.\,I.~Suriajaya, and C.~Turnage-Butterbaugh, ``Montgomery's pair correlation conjecture and its consequences for the distribution of zeros,'' \emph{Acta Arith.}~\textbf{214} (2024), 357--376.

\bibitem{BHB13} H.\,M.~Bui and D.\,R.~Heath-Brown, ``On simple zeros of the Riemann zeta-function,'' \emph{Bull.\ London Math.\ Soc.}~\textbf{45} (2013), 953--961.

\bibitem{CL55} E.\,A.~Coddington and N.~Levinson, \emph{Theory of Ordinary Differential Equations}, McGraw-Hill, 1955.

\bibitem{CM22} A.~Connes and H.~Moscovici, ``Prolate spheroidal operator and zeta,'' \emph{Proc.\ Natl.\ Acad.\ Sci.\ USA}~\textbf{119} (2022), No.\,22, e2123174119.

\bibitem{Con99} A.~Connes, ``Trace formula in noncommutative geometry and the zeros of the Riemann zeta function,'' \emph{Selecta Math.\ (N.S.)}~\textbf{5} (1999), 29--106.

\bibitem{CSV88} G.~Csordas, W.~Smith, and R.\,S.~Varga, ``Lehmer pairs of zeros, the de~Bruijn--Newman constant $\Lambda$, and the Riemann hypothesis,'' \emph{Constr.\ Approx.}~\textbf{10} (1994), 107--129.

\bibitem{dB50} N.\,G.~de~Bruijn, ``The roots of trigonometric integrals,'' \emph{Duke Math.\ J.}~\textbf{17} (1950), 197--226.

\bibitem{dB68} L.~de~Branges, \emph{Hilbert Spaces of Entire Functions}, Prentice-Hall, 1968.

\bibitem{EMN76} R.\,S.~Ellis, J.\,L.~Monroe, and C.\,M.~Newman, ``The GHS and other correlation inequalities for a class of even ferromagnets,'' \emph{Commun.\ Math.\ Phys.}~\textbf{46} (1976), 167--182.

\bibitem{EN78} R.\,S.~Ellis and C.\,M.~Newman, ``Necessary and sufficient conditions for the GHS inequality with applications to analysis and probability,'' \emph{Trans.\ Amer.\ Math.\ Soc.}~\textbf{237} (1978), 83--99.

\bibitem{GM78} P.\,X.~Gallagher and J.\,H.~Mueller, ``Primes and zeros in short intervals,'' \emph{J.\ Reine Angew.\ Math.}~\textbf{303/304} (1978), 205--220.

\bibitem{Kap22a} V.\,V.~Kapustin, ``The zero set of the Riemann zeta function as the point spectrum of an operator,'' \emph{Algebra i Analiz}~\textbf{33} (2021), No.\,4, 107--124; English transl., \emph{St.\,Petersburg Math.\,J.}~\textbf{33} (2022), 661--673.

\bibitem{Kap22b} V.\,V.~Kapustin, ``Five Hilbert space models related to the Riemann zeta function,'' \emph{J.\ Math.\ Sci.}~\textbf{268} (2022), 791--799.

\bibitem{Mon73} H.\,L.~Montgomery, ``The pair correlation of zeros of the zeta function,'' in \emph{Analytic Number Theory}, Proc.\ Sympos.\ Pure Math.~\textbf{24}, AMS, 1973, pp.\,181--193.

\bibitem{New91} C.\,M.~Newman, ``The GHS inequality and the Riemann hypothesis,'' \emph{Constr.\ Approx.}~\textbf{7} (1991), 389--399.

\bibitem{Odl} A.\,M.~Odlyzko, ``Correspondence about the origins of the Hilbert--P\'olya conjecture,'' available at \url{www-users.cse.umn.edu/~odlyzko/polya/}.

\bibitem{PT21} D.\,J.~Platt and T.\,S.~Trudgian, ``The Riemann hypothesis is true up to $3\times 10^{12}$,'' \emph{Bull.\ London Math.\ Soc.}~\textbf{53} (2021), 792--797.

\bibitem{RT20} B.~Rodgers and T.~Tao, ``The de~Bruijn--Newman constant is non-negative,'' \emph{Forum Math.\ Pi}~\textbf{8} (2020), e6.

\bibitem{Shr26c} S.~Shrestha, ``A proof of the Riemann Hypothesis via the GHS class and Lee--Yang property,'' 2026. Available at \url{https://github.com/Sudip-Shrestha0x0/ramanujan-project}.

\bibitem{Wei52} A.~Weil, ``Sur les `formules explicites' de la th\'eorie des nombres premiers,'' \emph{Meddelanden Lunds Univ.\ Mat.\ Sem.}, tome suppl\'ementaire (1952), 252--265.

\bibitem{Weyl10} H.~Weyl, ``\"Uber gew\"ohnliche Differentialgleichungen mit Singularit\"aten und die zugeh\"origen Entwicklungen willk\"urlicher Funktionen,'' \emph{Math.\ Ann.}~\textbf{68} (1910), 220--269.

\end{thebibliography}

\end{document}
